
\subsection{RQ2: Practical Analyzability}\label{sec:RQ2}

RQ2 asks whether observed bytecode artifacts can be analyzed by selected tools in a version-aware environment.  For each \texttt{.pyc} file found in RQ1, we run standard-library loading and disassembly, then selected decompilers when available.  Each tool result records both the status and the concrete failure reason, so failures caused by missing interpreters, unsupported bytecode versions, timeouts, and decompiler errors are not conflated.

\begin{table}[t]
  \centering
  \caption{Practical analyzability levels for observed bytecode artifacts. Replace placeholders with measured values.}
  \label{tab:tool-analyzability}
  \begin{tabular}{lrrrrr}
    \toprule
    Version & \texttt{.pyc} files & Loadable & Disassemblable & Partial decompile & Full decompile \\
    \midrule
    CPython 3.8 & TBD & TBD & TBD & TBD & TBD \\
    CPython 3.9 & TBD & TBD & TBD & TBD & TBD \\
    CPython 3.10 & TBD & TBD & TBD & TBD & TBD \\
    CPython 3.12 & TBD & TBD & TBD & TBD & TBD \\
    \bottomrule
  \end{tabular}
\end{table}

The key interpretation is practical rather than theoretical.  If bytecode becomes analyzable after the matching CPython interpreter is installed, the earlier failure is an environment limitation.  If selected tools still fail in a matched environment, the result indicates a current tooling limitation.  RQ2 therefore reports what the selected analysis stack can inspect today, not a proof that the bytecode is or is not recoverable by all possible tools.
